% ભાગ: ગણો-વિધેયો-સંબંધો
% પ્રકરણ: અંકગણિતીકરણ
% વિભાગ: Nથી Z

\documentclass[../../../include/open-logic-section]{subfiles}

\begin{document}

\olfileid{sfr}{arith}{int}
\olsection{$\Nat$થી $\Int$ સુધી}

અહીં બે મૂળભૂત અવલોકનો છે:
  \begin{enumerate}
    \item દરેક પૂર્ણાંકને $n - m$ સ્વરૂપે લખી શકાય છે, જ્યાં $n, m \in\Nat$.
    \item $n - m$ પદાવલિમાં રહેલી માહિતીને ક્રમયુક્ત જોડ $\tuple{n, m}$ વડે પણ બરાબર નિરૂપિત કરી શકાય છે.
  \end{enumerate}
પ્રાકૃતિક સંખ્યાઓની ક્રમયુક્ત જોડો $\Nat^2$ના !!{element}s છે તે આપણે
પહેલેથી જાણીએ છીએ. અને આપણે $\Nat$ને સમજીએ છીએ એમ ધારી રહ્યા છીએ.
તેથી આ બંને અવલોકનો પરથી એક સહજ સૂચન મળે છે: \emph{પૂર્ણાંકોને
પ્રાકૃતિક સંખ્યાઓની ક્રમયુક્ત જોડો તરીકે ગણીએ}.

હકીકતમાં આ સૂચન અતિસરળ છે. નિશ્ચિતપણે આપણે $0- 2 = 4 - 6$ થાય
એવું ઇચ્છીએ છીએ. પરંતુ દેખીતી રીતે $\tuple{0, 2 }\neq \tuple{4, 6}$.
આથી $\Nat^2$ એ જ પૂર્ણાંકોનો ગણ છે એવું સીધું કહી શકાતું નથી.

આ સમસ્યાનું સામાન્યકરણ કરતાં આપણને નીચેનું જોઈએ છે:
  $$a - b = c - d \text{ ત્યારે અને માત્ર ત્યારે જ જ્યારે }a + d = c + b$$
(પૂર્ણાંકોનો વ્યવહાર આવો જ \emph{હોવો જોઈએ} તે સ્પષ્ટ છે: બંને બાજુ
$b$ અને $d$ ઉમેરો.) આ વ્યવહાર સુનિશ્ચિત કરવાની સરળ રીત એટલે
ક્રમયુક્ત જોડો વચ્ચે સામ્ય સંબંધ $\Intequiv$ની નીચે મુજબ વ્યાખ્યા કરવી:
$$\tuple{ a, b } \Intequiv \tuple{c, d}\text{ ત્યારે અને માત્ર ત્યારે જ જ્યારે }a + d = c + b$$
હવે આપણે બતાવવું પડશે કે આ સામ્ય સંબંધ છે.
\begin{prop} $\Intequiv$ એ સામ્ય સંબંધ છે.
\end{prop}
  \begin{proof}
  આપણે બતાવવું પડશે કે $\Intequiv$ સ્વવાચક, સંમિત અને પરંપરિત છે.

  \emph{સ્વવાચકતા:} $a + b = b + a$ હોવાથી દેખીતી રીતે
  $\tuple{a, b} \Intequiv \tuple{a, b}$.

  \emph{સંમિતતા:} ધારો કે $\tuple{a, b} \Intequiv \tuple{c, d}$,
  એટલે કે $a + d = c + b$. ત્યારે $c + b = a + d$, તેથી
  $\tuple{c, d} \Intequiv \tuple{a, b}$.

  \emph{પરંપરિતતા:} ધારો કે
  $\tuple{a, b} \Intequiv \tuple{c, d}\Intequiv \tuple{m, n}$.
  તેથી $a + d = c + b$ અને $c + n = m + d$. આથી
  $a + d + c + n = c + b + m + d$, અને તેથી $a + n = m + b$.
  પરિણામે $\tuple{a, b } \Intequiv \tuple{m, n}$.
\end{proof}

હવે આ સામ્ય સંબંધ વડે આપણે સામ્ય વર્ગો લઈ શકીએ છીએ:

\begin{defn}
    પૂર્ણાંકો એ $\Intequiv$ હેઠળ પ્રાકૃતિક સંખ્યાઓની ક્રમયુક્ત
    જોડોના સામ્ય વર્ગો છે; એટલે કે,
    $\Int = \equivclass{\Nat^2}{\Intequiv}$.
\end{defn}

આ નિર્ધારક વ્યાખ્યા પ્રત્યે જુદી જુદી \emph{તાત્ત્વિક} પ્રતિક્રિયાઓ
હોઈ શકે. જોકે એ પ્રતિક્રિયાઓ વિચારીએ તે પહેલાં કેટલીક પ્રાવિધિક
વિગતો આગળ વધારી લેવી ઉપયોગી છે.

પૂર્ણાંકો શું છે તે કહ્યા પછી તેમના પરનાં મૂળભૂત વિધેયો અને સંબંધોની
વ્યાખ્યા કરવી પડશે. $\equivrep{m, n}{\Intequiv}$ વડે આપણે
$\Intequiv$ હેઠળનો એવો સામ્ય વર્ગ દર્શાવીશું જેમાં $\tuple{m, n}$
!!a{element} છે.\footnote{નોંધ: \olref[sfr][rel][eqv]{def:equivalenceclass}માં
રજૂ કરેલી સંજ્ઞા વાપરીએ તો આ જ વસ્તુ માટે
$\equivrep{\tuple{m,n}}{\Intequiv}$ લખાત. પરંતુ તે વાંચવામાં થોડું
મુશ્કેલ છે.} એટલે કે:
  $$\equivrep{m, n}{\Intequiv} = \Setabs{\tuple{a, b} \in \Nat^2}{\tuple{a, b}\Intequiv \tuple{m, n}}$$
હવે આપણે નીચેની વ્યાખ્યાઓ આપીએ છીએ:
  \begin{align*}
    \equivrep{a, b}{\Intequiv} + \equivrep{c, d}{\Intequiv} &= \equivrep{a + c, b + d}{\Intequiv}\\
    \equivrep{a, b}{\Intequiv} \times \equivrep{c, d}{\Intequiv} &= \equivrep{a c + b  d, a  d + b c}{\Intequiv}\\
    \equivrep{a, b}{\Intequiv} \leq \equivrep{c, d}{\Intequiv} &\text{ ત્યારે અને માત્ર ત્યારે જ જ્યારે }a + d \leq b + c
  \end{align*}
(સામાન્ય પ્રથા મુજબ સ્વયંસિદ્ધ સૂત્રો સરળતાથી વાંચી શકાય તે માટે
હું `$(a \times b)$'ને બદલે `$ab$' લખું છું.) હવે આ વ્યાખ્યાઓનો
વ્યવહાર જેવો \emph{હોવો જોઈએ} તેવો જ છે તેની ખાતરી કરવી પડશે.
તેનો અર્થ વિગતે લખીને તપાસ કરવી ઘણી લાંબી પ્રક્રિયા છે; તેથી વિગતો
\olref[check]{sec}માં રાખી છે. પરંતુ ટૂંકી વાત એ છે કે બધું કામ કરે છે!

છેલ્લી એક વાત બાકી છે. આપણે પ્રાકૃતિક સંખ્યાઓનો ઉપયોગ કરીને
પૂર્ણાંકોની રચના કરી છે. પરંતુ તેનો અર્થ એ થશે કે પ્રાકૃતિક સંખ્યાઓ
\emph{પોતે પૂર્ણાંકો નથી}. આનું તાત્ત્વિક મહત્ત્વ આપણે
\olref[ref]{sec}માં ફરી વિચારીશું. જોકે નિર્વિવાદ પ્રાવિધિક સ્તરે
પ્રાકૃતિક સંખ્યાઓને પૂર્ણાંકો \emph{તરીકે} ગણવાની કોઈ રીત જોઈએ.
વિચાર બહુ સરળ છે: દરેક $n \in \Nat$ માટે આપણે માત્ર એમ ઠરાવીએ કે
$n_\Int = \equivrep{n, 0}{\Intequiv}$. આ વ્યાખ્યા સુવ્યવસ્થિત છે,
એટલે કે કોઈ પણ $m, n \in \Nat$ માટે નીચેનું થાય છે, તેની ખાતરી કરવી પડશે:
\begin{align*}
  (m + n)_\Int &= m_\Int + n_\Int\\
  (m \times n)_\Int &= m_\Int \times n_\Int\\
  m \leq n &\liff m_\Int \leq n_\Int
\end{align*}
પરંતુ આ બધું ઘણું સીધું છે. દાખલા તરીકે, આમાંની બીજી સમાનતા
બતાવવા પ્રાકૃતિક સંખ્યાઓના વ્યવહારનો મુક્તપણે ઉપયોગ કરીને નીચે મુજબ
દલીલ કરી શકીએ છીએ:
%\begin{proof}
%  પ્રાકૃતિક સંખ્યાઓના વ્યવહારનો ઉપયોગ કરતાં દેખીતી રીતે:
  \begin{align*}
%    (m + n)_\Int  = \equivrep{m + n, 0}{\Intequiv} = \equivrep{m + n, 0 + 0}{\Intequiv} = \equivrep{m, 0}{\Intequiv} + \equivrep{n, 0}{\Intequiv} = m_\Int + n_\Int\\
    (m \times n)_\Int  &= \equivrep{m \times n, 0}{\Intequiv} \\
    &= \equivrep{m \times n + 0 \times 0, m \times 0 + 0 \times n}{\Intequiv} \\
    &= \equivrep{m, 0}{\Intequiv} \times \equivrep{n, 0}{\Intequiv} \\
    &= m_\Int \times n_\Int
  \end{align*}
%\end{proof}
બીજી બે શરતો સાચી છે તેની ખાતરી કરવાનું સ્વાધ્યાય તરીકે રાખીએ છીએ.
\begin{prob}
  કોઈ પણ $m, n \in \Nat$ માટે $(m + n)_\Int = m_\Int + n_\Int$
  અને $m \leq n \liff m_\Int \leq n_\Int$ છે તે બતાવો.
\end{prob}
\end{document}
